\subsection{A geometric realization of the non-self-dual periods}

We now construct periods for the specific integral datum of
Theorems F.2.5--F.2.6.  The source below is a family of genus-two
curves, but the rank-six system is an extension of its primitive
exterior square.  It does not have the filtration excluded in
Theorem F.3.4.

\paragraph{Theorem F.3.5 (Proved: the geometric rank-five quotient).}
Put $\lambda_*=-256/3125$ and
$U=\mathbb P^1\setminus\{0,\infty,\lambda_*\}$.  The smooth projective
models of
\[
 C_\lambda:\quad y^2=x^5+\lambda^3x+\lambda^4
 \tag{F.7}
\]
give a polarized rank-four rational cohomology system $\mathcal H$ on
$U$.  Its primitive exterior square is rationally isomorphic to
the rank-five quotient of the system with matrices $A,B$ in
Theorem F.2.5 by its invariant first coordinate line.
On the orbifold $\mathcal B'$ obtained by taking roots of orders
five at infinity and eight at zero, the extended Hodge bundle is
\[
 \mathcal E=\mathcal T^3\oplus\mathcal T,\qquad
 \mathcal T=\mathcal O_{\mathcal B'}(2D_\infty-3D_0),\qquad
 \deg\mathcal T=\frac1{40}.
 \tag{F.8}
\]
Here $D_\infty,D_0$ are the reduced root divisors; the extension at
$\lambda_*$ is the nilpotent canonical extension.

\begin{proof}
The discriminant of the degree-five polynomial in (F.7) is exactly
$\lambda^{15}(3125\lambda+256)$.  At $\lambda_*$ its unique double
root is $r=64/625$ and the coefficient of $(x-r)^2$ is
$10r^3\ne0$.  Its other roots are simple.  Thus the fibre has one
ordinary node and irreducible genus-one normalization.  The local
model is a smoothing of a nonseparating node; the vanishing circle
is a primitive nonseparating homology class.  Cutting and regluing
an annular neighborhood proves that its monodromy is a single
Dehn twist, and hence a primitive rank-one transvection on $H^1$.
In particular its logarithm $N_H$ has square zero and rank one.

At zero, the substitutions $\lambda=s^8$, $x=s^6X$, $y=s^{15}Y$
give
\[
 Y^2=X^5+X+s^2.
\]
This is smooth at $s=0$.  The generator $s\mapsto\zeta_8s$ at
fixed $(x,y)$ acts by $(X,Y)\mapsto(iX,\zeta_8Y)$, so its
eigenvalues on the two holomorphic differentials are
$\zeta_8,\zeta_8^3$.  The full rational characteristic polynomial
is $g(T)=T^4+1$.
At infinity put $\lambda=s^{-5}$, $x=s^{-4}X$, $y=s^{-10}Y$.
The resulting smooth family is
\[
 Y^2=X^5+sX+1.
\]
The deck generator acts by $(X,Y)\mapsto(\zeta_5^4X,Y)$,
with holomorphic eigenvalues $\zeta_5^4,\zeta_5^3$.  Its full
characteristic polynomial is $f(T)=T^4+T^3+T^2+T+1$.
Inverses have the same two rational characteristic polynomials.

We identify the representation, rather than merely its spectra.
Choose positive peripheral generators $T_1,T_2$ at infinity and
zero, in that order, with $T_1T_2$ inverse to nodal monodromy.
Set $R=T_1^{-1}$ and $S=T_2$.  Then $R-S$ has rank one and their
characteristic polynomials are $f,g$.  The following elementary
rank-one lemma fixes their simultaneous rational conjugacy class.
If $u$ spans $\operatorname{im}(R-S)$, it is cyclic for $R$:
the span of its iterates is stable under both $R$ and $S$, and
the irreducibility of $f$ makes that span all of $\mathbb Q^4$.
In the rational basis $u,Ru,R^2u,R^3u$, the matrix of $R$ is
the companion matrix of $f$, and $S=R-u\varphi$ for a covector
$\varphi$.  The determinant identity
\[
 g(T)=f(T)+\varphi\operatorname{adj}(TI-R)u
 \tag{F.9}
\]
determines $\varphi$ uniquely: the four coordinates of
$\operatorname{adj}(TI-R)u$ are a basis of the polynomials of
degree at most three, as is seen from their successively distinct
leading degrees in the cyclic basis.  It follows that any two
such pairs are simultaneously rationally conjugate.  Apply the
lemma to the explicit companion pair
\[
 C_f=\begin{pmatrix}
 0&0&0&-1\\1&0&0&-1\\0&1&0&-1\\0&0&1&-1
 \end{pmatrix},\qquad
 C_g=\begin{pmatrix}
 0&0&0&-1\\1&0&0&0\\0&1&0&0\\0&0&1&0
 \end{pmatrix}.
 \tag{F.10}
\]
Their difference has rank one and their characteristic polynomials
are $f,g$.  Thus $(T_1,T_2)$ is rationally conjugate to
$(C_f^{-1},C_g)$, with the chosen peripheral conventions.

Use the lexicographic basis in $\bigwedge^2\mathbb Q^4$.
Its common invariant vector is $(1,1,1,1,1,1)^t$.
Let $Q$ be the identity matrix of size six with first column
replaced by this vector.  It is unimodular.  Direct multiplication
gives the exact identities
\[
 \left(Q^{-1}\bigwedge^2(C_f^{-1})Q\right)_{2:6,2:6}
   =A_{2:6,2:6},\qquad
 \left(Q^{-1}\bigwedge^2(C_g)Q\right)_{2:6,2:6}
   =B_{2:6,2:6}.
 \tag{F.11}
\]
They are checked over $\mathbb Z$ in
\texttt{fourfolds\_periods.py}, not inferred from eigenvalues.
The invariant line is the inverse of the polarization, up to a
nonzero rational scalar.  Projection along that line identifies
the quotient with the primitive exterior square $\mathcal W$.

The differentials $\omega_1=dx/y$ and $\omega_2=x\,dx/y$ span
the extended Hodge bundle everywhere.  At the node they are
independent dualizing differentials: $\omega_1$ has nonzero
opposite residues on the normalization and
$(x-r)\omega_1$ is a nonzero holomorphic differential there.
Here the extension assertion can be checked on the nodal
annulus, without invoking an additional extension theorem.
In coordinates $uv=t$ a relative dualizing differential has
a convergent Laurent expansion
$\sum_n a_n(t)u^n\,du/u$.  Integration on the vanishing circle
is $2\pi i a_0(t)$.  Integration across the annulus has the
single logarithmic term $a_0(t)\log t$; integration of every
nonconstant Laurent monomial and of the complementary fixed
arcs gives a holomorphic remainder.  Thus the period vector
has the form $\exp(sN_H)$ times a holomorphic vector,
$s=\log t/(2\pi i)$.  Untwisting by $\exp(-sN_H)$ is exactly
the nilpotent canonical extension used here.  The two independent
dualizing limits remain a basis of its Hodge subbundle.
Their isotropic pairing, zero away from zero, remains zero
at the limit by holomorphic continuation, so that limit is
Lagrangian.
At zero their expressions in the root coordinates are
$s^{-9}dX/Y$ and $s^{-3}X\,dX/Y$; at infinity they are
$s^6dX/Y$ and $s^2X\,dX/Y$.
Thus their divisor lines are
$\mathcal O(6D_\infty-9D_0)$ and
$\mathcal O(2D_\infty-3D_0)$, proving (F.8).
The nonzero residue of $\omega_1$ also shows
\[
 N_H\ell_\infty\ne0,
 \qquad \ell=\mathcal T^3\subset\mathcal E,
 \tag{F.12}
\]
where $\ell_\infty$ denotes its limit at the nodal cusp.
Indeed integration on the vanishing circle tends to
$2\pi i$ times that residue; the transvection sends the cohomology
class of a differential to that period times its vanishing
cohomology vector.
\end{proof}

\paragraph{Theorem F.3.6 (Proved: a global nonpolarized period family).}
The integral rank-six system with cohomology monodromies $A,B$
admits a holomorphic subbundle $\mathcal F$ of rank three over
$U$ satisfying $\mathcal F\cap\overline{\mathcal F}=0$.
It extends holomorphically on the finite root charts.  At the
cusp an integral marking puts its periods in the form
\[
 Z(t)=\frac{\log t}{2\pi i}I_3+C(t),
 \tag{F.13}
\]
where $C$ is holomorphic at zero and bounded there.

\begin{proof}
Write the exact sequence supplied by (F.11) as
\[
 0\longrightarrow\mathbb Q e\longrightarrow\mathcal V
   \longrightarrow\mathcal W\longrightarrow0.
 \tag{F.14}
\]
Here $e$ is the first coordinate vector.  The nilpotent canonical
extension at the cusp and the ordinary extensions on finite root
charts preserve this sequence.  This follows directly by using
the frames $\exp(-sN)$: each residue has only eigenvalue zero,
and the same exponential construction respects the inclusion and
the quotient.  At finite roots the statement is ordinary exactness
of bundles with trivialized local monodromy.  Denote these extended
bundles by $\mathcal H_D,\mathcal V_D,\mathcal W_D$.

Let $\omega\in\bigwedge^2\mathcal H_D$ be the real, flat
nondegenerate bivector inverse to the polarization, and let
$p:\bigwedge^2\mathcal H_D\to\mathcal W_D$ be primitive projection.
Define the holomorphic rank-three bundle
\[
 \mathcal P=p(\ell\wedge\mathcal H_D).
 \tag{F.15}
\]
Projection is injective on $\ell\wedge\mathcal H_D$: every
nonzero bivector there is decomposable whereas $\omega$ is not.
This argument applies also at the cusp.  The polarization gives
$\mathcal H_D/\mathcal E\simeq\mathcal E^*$, including there,
since $\mathcal E$ is Lagrangian.  Hence
\[
 0\longrightarrow\mathcal T^4\longrightarrow\mathcal P
 \longrightarrow\mathcal O\oplus\mathcal T^2\longrightarrow0.
 \tag{F.16}
\]
After dualizing, the end terms are
$\mathcal O\oplus\mathcal T^{-2}$ and $\mathcal T^{-4}$.

We compute the obstruction to lifting $\mathcal P$ in (F.14).
For the coarse projection $\rho:\mathcal B'\to\mathbb P^1$,
the elementary root-chart calculation of invariant sections gives
\[
 \rho_*\mathcal T^k=
 \mathcal O_{\mathbb P^1}
 \left(\left\lfloor\frac{2k}{5}\right\rfloor+
       \left\lfloor\frac{-3k}{8}\right\rfloor\right).
 \tag{F.17}
\]
Higher direct images vanish, because taking invariants under a
finite group in characteristic zero is exact.  For $k=0,-2,-4$
the displayed degrees are respectively $0,-1,-1$.  Their first
cohomology vanishes.  The dual of (F.16) therefore gives
$H^1(\mathcal B',\mathcal P^*)=0$.
The pullback of (F.14) to $\mathcal P$ consequently splits as
holomorphic bundles: choose a holomorphic lift
$\sigma:\mathcal P\to\mathcal V_D$ of its inclusion in
$\mathcal W_D$.  This is an actual extension-class calculation
on a compact orbifold, not an approximation or an Oka principle.
To specify the choice concretely, take a finite root-coordinate
cover with local bundle lifts $\sigma_i$ and write their
differences as $\xi_{ij}e$, where
$\xi_{ij}$ is a Cech cocycle in $\mathcal P^*$.
Choose any solution $b_i-b_j=\xi_{ij}$ of this explicitly
vanishing cohomology class and glue the maps
$\sigma_i-b_ie$.  Every such solution is allowed in the
construction that follows.

There is also a canonical holomorphic functional
$\alpha:\mathcal P\to\mathcal O$.
For $v\in\mathcal P$, its unique preimage in
$\ell\wedge\mathcal H_D$ is
\[
 j(v)=v+\alpha(v)\omega
 \tag{F.18}
\]
under the primitive decomposition.
At a smooth point, the positive Hodge metric implies that
$\ell$ and $\bar\ell$ form a nondegenerate symplectic pair.
Their common wedge line is $\ell\wedge\bar\ell$.
The bivector $\omega$ does not belong to
$\ell\wedge\mathcal H+\bar\ell\wedge\mathcal H$, since its
component on the two-dimensional symplectic orthogonal complement
is nonzero.  Projection $p$ therefore preserves the intersection,
and
\[
 \mathcal P\cap\bar{\mathcal P}
      =p(\ell\wedge\bar\ell).
 \tag{F.19}
\]
This is the complexification of a real line.  If $u$ spans
$\ell$, then $i u\wedge\bar u$ is real and its coefficient on
$\omega$ is nonzero, by positivity.  Thus $\alpha$ is real and
nonzero on that real intersection line.  There is a unique smooth
real section $n$ of it such that $\alpha(n)=1$.

Since the quotient of $\sigma(n)$ is real, there is a unique
real smooth function $h$ on $U$ such that
\[
 \sigma(n)-\overline{\sigma(n)}=2ih e.
 \tag{F.20}
\]
For $c\in\mathbb R$ set
\[
 \mathcal F_c=(\sigma+ic\alpha e)(\mathcal P).
 \tag{F.21}
\]
This is a holomorphic rank-three subbundle projecting
isomorphically to $\mathcal P$.  An intersection with its complex
conjugate must project to (F.19).  On its real generator the two
lifts differ by $2i(h+c)e$.  It follows, including arbitrary
complex multiples of that generator, that
\[
 \mathcal F_c\cap\overline{\mathcal F_c}=0
       \quad\Longleftrightarrow\quad h+c\ne0.
 \tag{F.22}
\]

We verify the required cusp property before choosing $c$.
In a complex symplectic basis $a,b,d_1,d_2$ of the limiting
$\mathcal H_D$, relation (F.12) permits us to arrange
$N_Hb=a$, $N_Ha=N_Hd_1=N_Hd_2=0$ and $\ell_\infty=\mathbb Cb$.
Then
\[
 \mathcal P_\infty=
 \langle k=p(b\wedge a),\ b\wedge d_1,\ b\wedge d_2\rangle,
 \quad
 N_W(b\wedge d_i)=a\wedge d_i,
 \quad N_Wk=0.
 \tag{F.23}
\]
The line generated by $k$ is
$\ker N_W/\operatorname{im}N_W$; all other kernel directions
come from the displayed image.  The exact matrices give
$N_V^2=0$, $\operatorname{rank}N_V=3$ and
$\operatorname{rank}N_W=2$.
Every lift of an image vector $a\wedge d_i$ is killed by $N_V$,
because that vector has a lift in $\operatorname{im}N_V$ and
$N_Ve=0$.  Rank three now forces $N_V$ to take any lift of $k$
to a nonzero multiple of $e$.  Therefore
$N_V:\mathcal F_{c,\infty}\to\operatorname{im}N_V$ is an
isomorphism for every $c$.  In particular its domain is a
complement to $\ker N_V$.

The primitive index-one cusp has a real integral basis
$a_i,b_i$ with $N_Vb_i=a_i$, by Proposition F.2.1, also for the
dual or inverse convention.  In this basis the untwisted
subbundle near zero is the graph $b+C_c(t)a$, with $C_c$
holomorphic and bounded.  Restoring the flat frame by
$\exp(sN_V)$, $s=\log t/(2\pi i)$, gives (F.13).
Thus $\operatorname{Im}Z=\operatorname{Im}s\,I_3+
\operatorname{Im}C_c$ is nonsingular on a sufficiently small
punctured disc.  Apply this first with $c=0$.
Equation (F.22) shows that $h$ is nonzero and has constant sign
on that connected punctured disc.

At the two finite root points, all ingredients in (F.19)--(F.20)
are smooth and $\alpha(n)=1$ remains valid, so $h$ extends and
is bounded.  On the compact complement of the cusp disc it has
a finite bound $K$.  Choose $c$ with the same sign as $h$ near
the cusp by the exact rule
$c=\operatorname{sign}(h(t_0))(K+2)$, where $t_0$ is any point
of that punctured disc and $K$ is the supremum of $|h|$ on
the compact complement.  Then $h+c$ is nonzero both on that
disc and on its complement.  Equations (F.21)--(F.22) prove
global real transversality, and the preceding cusp calculation
proves the bounded correction for this same $c$.

Finally, the fourth power of the order-eight action on
$\mathcal H$ is $-I$.  It acts trivially on
$\mathcal W,\mathcal P,\alpha$ and $\mathcal V$.
Hence these bundles, the lift and the constructed periods descend
from $\mathcal B'$ to the orbifold $\mathcal B$ with orders five
and four.  Indeed any equivariant bundle map between bundles with
trivial action of the extra order-two stabilizer descends on its
root chart.  This gives exactly the finite orders required for
the rank-six datum.
\end{proof}

\paragraph{Theorem F.3.7 (Proved: compact holomorphic filling).}
The smooth eight-manifold of Theorem F.2.6 carries an integrable
complex structure.  There is a proper holomorphic map to
$\mathbb P^1$ whose regular fibres are complex three-tori,
whose two finite reduced fibres are the free affine quotients
of orders five and four specified in Theorem F.2.5, and whose
remaining fibre is the index-one Kuhn degeneration.

\begin{proof}
We first fix the dual convention.  The preceding construction
has cohomology monodromies $A,B$; Theorem F.2.5 uses homology
monodromies $A^t,B^t$.  All three special values are real.
Choose a positive real basepoint and paths along the positive
real circle to infinity and to zero.  Complex conjugation fixes
these paths and reverses the respective peripheral meridians.
Pull back the preceding period family by conjugation of the
base and conjugate its complex fibres.  Its Hodge subbundle is
holomorphic, with local expression
$\overline{\mathcal F(\bar\lambda)}$, and its cohomology
monodromies are $A^{-1},B^{-1}$.  Its homology monodromies are
therefore exactly $A^t,B^t$.  The third monodromy is fixed by
the product relation, so this also fixes its cusp marking.
The primitive index-one condition and the bounded holomorphic
cusp correction are preserved, after replacing the real
angular basis by its negative if necessary.

For a holomorphic Hodge subbundle $\mathcal F$ with
$\mathcal F\cap\bar{\mathcal F}=0$, restriction of integral
homology functionals embeds the homology lattice discretely in
$\mathcal F^*$.  Real transversality makes its real span the
whole real vector space.  The quotient
$\mathcal F^*/\mathcal V_{\mathbb Z}^*$ is consequently a
holomorphic family of compact complex three-tori over $U$,
with its zero section.  This also proves local properness,
by a smoothly varying fundamental parallelepiped.

On the root disc at the finite point of order $m_i$, the family
extends smoothly and its deck map has linear part $M_i$.
The invariant vector $v_i$ of Theorem F.2.5 gives a holomorphic
torsion section $v_i/m_i$.  Use the diagonal affine deck action
\[
 (z,s)\longmapsto(M_i z+v_i/m_i,\zeta_{m_i}s).
 \tag{F.24}
\]
Here $M_i z$ denotes the induced complex-linear map
$\mathcal F_s^*\to\mathcal F_{\zeta_{m_i}s}^*$ from the marked
lattice action, not multiplication of a complex three-vector
by a six-by-six matrix.  Likewise $v_i/m_i$ means its evaluation
as a point of the target torus.
It has order $m_i$, since the $m_i$th translation is integral.
Every nonidentity power is fixed-point-free on the central
torus: the invariant integral covector $\ell$ takes $v_i$ to
$1$ or $-1$, so a fixed point would require the noninteger
$\pm j/m_i$ to vanish in $\mathbb R/\mathbb Z$.
Away from the central fibre the base rotation is already free.
The quotient is therefore smooth and proper over the coarse
disc, with precisely the required multiple fibre.

Its marked punctured family is holomorphically identified with
the original torus family by logarithmic translation.  On the
root cover use $r(s)=(\log s/(2\pi i))v_i$.  Then
$r(\zeta_{m_i}s)-M_i r(s)=v_i/m_i$, modulo the integral lattice;
the change of branch of $\log s$ is itself an integral vector.
Thus translation by $r$ conjugates the linear deck action to
(F.24).  This fixes the affine clutching, including its sign.

At the cusp use the integral marking of (F.13), the fan (E.4)
and the action (E.5) with $C_3=C$.
Its maximal charts are $\mathbb C^4$ with map
$t=z_0z_1z_2z_3$, and its transition maps are the toric changes
of coordinates multiplied by the holomorphic units
$\exp(2\pi i C(t)\lambda)$.
The proof of Theorem E.1.2 used only boundedness of $C$,
integrality of the standard cusp, and unimodularity of that
fan.  It therefore proves smoothness, proper discontinuity
and properness for this model without a symmetry condition
on $C$.  Glue by the displayed integral marking and use no
additional cusp translation.

The three punctured overlap maps are holomorphic isomorphisms
with the same marked torus family, so their cocycle condition
is literal.  Proper local models and a compact complementary
base give a Hausdorff compact total space.  Every chart is
smooth.  The affine vectors, cusp fan and clutching are exactly
those in Theorems F.2.5--F.2.6.  Varying their bounded corrections
and varying the fibre complex structures along a smooth path
does not change these marked smooth fillings: on the finite
root models the underlying real torus family is the fixed
local-system bundle; on the cusp use the stratum-preserving
parameter vector fields and uniform properness estimate of
Theorem E.1.3.  Thus the resulting underlying smooth manifold
is the one whose integral topology was computed there.
\end{proof}
